\documentclass[10pt,a4paper]{amsart}
\usepackage{amsmath,amssymb,amsthm,graphicx}
\usepackage{verbatim}
\newenvironment{reference}{\comment}{\endcomment}
%\newenvironment{reference}{}{}
\newenvironment{slogan}{\comment}{\endcomment}
\newenvironment{history}{\comment}{\endcomment}

% For commutative diagrams we use Xy-pic
\usepackage[all]{xy}

% We use 2cell for 2-commutative diagrams.
\xyoption{2cell}
% Theorem environments.
%
\theoremstyle{plain}
\newtheorem{theorem}[subsection]{Theorem}
\newtheorem{proposition}[subsection]{Proposition}
\newtheorem{lemma}[subsection]{Lemma}

\theoremstyle{definition}
\newtheorem{definition}[subsection]{Definition}
\newtheorem{example}[subsection]{Example}
\newtheorem{exercise}[subsection]{Exercise}
\newtheorem{situation}[subsection]{Situation}

\theoremstyle{remark}
\newtheorem{remark}[subsection]{Remark}
\newtheorem{remarks}[subsection]{Remarks}

\numberwithin{equation}{subsection}

% Macros
%
\def\lim{\mathop{\mathrm{lim}}\nolimits}
\def\colim{\mathop{\mathrm{colim}}\nolimits}
\def\Spec{\mathop{\mathrm{Spec}}}
\def\Hom{\mathop{\mathrm{Hom}}\nolimits}
\def\Ext{\mathop{\mathrm{Ext}}\nolimits}
\def\SheafHom{\mathop{\mathcal{H}\!\mathit{om}}\nolimits}
\def\SheafExt{\mathop{\mathcal{E}\!\mathit{xt}}\nolimits}
\def\Sch{\mathit{Sch}}
\def\Mor{\mathop{\mathrm{Mor}}\nolimits}
\def\Ob{\mathop{\mathrm{Ob}}\nolimits}
\def\Sh{\mathop{\mathit{Sh}}\nolimits}
\def\NL{\mathop{N\!L}\nolimits}
\def\CH{\mathop{\mathrm{CH}}\nolimits}
\def\proetale{{pro\text{-}\acute{e}tale}}
\def\etale{{\acute{e}tale}}
\def\QCoh{\mathit{QCoh}}
\def\Ker{\mathop{\mathrm{Ker}}}
\def\Im{\mathop{\mathrm{Im}}}
\def\Coker{\mathop{\mathrm{Coker}}}
\def\Coim{\mathop{\mathrm{Coim}}}

% Boxtimes
%
\DeclareMathSymbol{\boxtimes}{\mathbin}{AMSa}{"02}

%
% Macros for moduli stacks/spaces
%
\def\QCohstack{\mathcal{QC}\!\mathit{oh}}
\def\Cohstack{\mathcal{C}\!\mathit{oh}}
\def\Spacesstack{\mathcal{S}\!\mathit{paces}}
\def\Quotfunctor{\mathrm{Quot}}
\def\Hilbfunctor{\mathrm{Hilb}}
\def\Curvesstack{\mathcal{C}\!\mathit{urves}}
\def\Polarizedstack{\mathcal{P}\!\mathit{olarized}}
\def\Complexesstack{\mathcal{C}\!\mathit{omplexes}}
% \Pic is the operator that assigns to X its picard group, usage \Pic(X)
% \Picardstack_{X/B} denotes the Picard stack of X over B
% \Picardfunctor_{X/B} denotes the Picard functor of X over B
\def\Pic{\mathop{\mathrm{Pic}}\nolimits}
\def\Picardstack{\mathcal{P}\!\mathit{ic}}
\def\Picardfunctor{\mathrm{Pic}}
\def\Deformationcategory{\mathcal{D}\!\mathit{ef}}
\usepackage{cleveref}
\crefname{originalreference}{original reference}{original references}
\Crefname{originalreference}{Original reference}{Original references}
\makeatletter
\newlabel{item-construction-associated-cofibered-groupoid}{{06GN}{source}{}{sourcepointer-1}{}}
\newlabel{item-construction-associated-cofibered-groupoid@cref}{{[originalreference][0][]06GN}{[1][0][]source}{}{}{}}
\makeatother
\crefname{originalreference}{original reference}{original references}
\Crefname{originalreference}{Original reference}{Original references}
\makeatletter
\newlabel{remarks-cofibered-groupoids}{{06GK}{source}{}{sourcepointer-2}{}}
\newlabel{remarks-cofibered-groupoids@cref}{{[originalreference][0][]06GK}{[1][0][]source}{}{}{}}
\makeatother
\title{The quotient of a smooth groupoid in functors whose objects and arrows satisfy (RS) also satisfies (RS)}
\author{The Stacks Project Authors}
\date{stacks-project:formal-defos.tex:local-snapshot-20261012; GNU FDL 1.2 or later}
\begin{document}
\maketitle
\noindent\textit{Editorial note.} All delivered mathematical text is the author's own, copied from the pinned Stacks Project chapter file formal-defos.tex; no statement or proof step was rewritten, reordered or removed.\par\medskip
\noindent\textit{Editorial note.} The counted proof is the author's complete argument: the translation of a diagram in the quotient into ring and element data, the construction of the candidate fibre product through smoothness of the source map and the (RS) property for the objects, and the verification of the universal property using the (RS) property for the arrows.\par\medskip
\noindent\textit{Editorial note.} The label lemma-smooth-RS-groupoid-in-functors-quotient names the counted statement, and the proof environment immediately following it in the source is delivered verbatim as its proof.\par\medskip
\noindent\textit{Editorial note.} The counted statement is carried by the author's own lemma environment in the chapter source.\par\medskip
\noindent\textit{Editorial note.} Pre-open gate: python operations/agent-authoring/tools/precard\_gate.py --source formal-defos.tex --label lemma-smooth-RS-groupoid-in-functors-quotient --start 6713 --end 6789 --local-id 28016 returned RESULT=OK (L28016 free, occupancy FREE, public interval clear). Candidate taken from sources/stacks-formal-defos/candidates-halfA.json.\par\medskip
\noindent\textit{Editorial note.} The statement and proof are copied from the pinned chapter file; the cited inputs are kept as external pointers with their canonical Stacks tags.\par\medskip
\noindent\textit{Allowed standard prerequisites (editorial).}\begin{itemize}
\item Smoothness interpretation for the source and target maps. Same chapter, the interpretation of smoothness of the morphism s : R -> U as the existence of lifts of elements along surjective ring maps. Supplies the element a-tilde over A\_2 lifting a along f\_2 and mapping to x\_2 under s.
\item The (RS) property for U and for R. Same chapter, the statement of the (RS) condition as the identification of the fibre product over A with the fibres over A\_1 and A\_2. Supplies the identifications U(A\_1) x\_\{U(A)\} U(A\_2) = U(A\_1 x\_A A\_2) and R(A\_1 x\_A A\_2) = R(A\_1) x\_\{R(A)\} R(A\_2).
\item Description of the quotient groupoid. Same chapter, the remark describing the objects and morphisms of the quotient [U/R] as pairs (A, x) and (f, a) with the compatibility conditions used in the proof. Supplies the translation of diagrams in [U/R] into the data f\_i, x\_i and a\_i used throughout.
\end{itemize}\medskip
\ifdefined\hypertarget\hypertarget{sourcepointer-1}{}\fi
\noindent\textit{Source reference.} Original reference 06GN: sources/stacks-formal-defos/pinned-original-source.tex, source lines 1432--1432. Prerequisite pointer invoked at proof line(s) 6730; it names a standard fact of the same chapter that the counted argument uses, and every step of the counted proof built on it is retained verbatim in the delivered body.\par\medskip
\ifdefined\hypertarget\hypertarget{sourcepointer-2}{}\fi
\noindent\textit{Source reference.} Original reference 06GK: sources/stacks-formal-defos/pinned-original-source.tex, source lines 1381--1381. Prerequisite pointer invoked at proof line(s) 6729; it names a standard fact of the same chapter that the counted argument uses, and every step of the counted proof built on it is retained verbatim in the delivered body.\par\medskip
\begin{lemma}
\label{lemma-smooth-RS-groupoid-in-functors-quotient}
Let $(U, R, s, t, c)$ be a smooth groupoid in functors on $\mathcal{C}_\Lambda$.
Assume $U$ and $R$ satisfy (RS). Then $[U/R]$ satisfies (RS).
\end{lemma}

\begin{proof}
Let
$$
\xymatrix{
                           &     (A_2, x_2) \ar[d]^{(f_2, a_2)} \\
(A_1, x_1) \ar[r]^{(f_1, a_1)} &     (A, x)
}
$$
be a diagram in $[U/R]$ such that $f_2: A_2 \to A$ is surjective. The
notation is as in
Remarks \ref{remarks-cofibered-groupoids}
(\ref{item-construction-associated-cofibered-groupoid}).
Hence $f_1: A_1 \to A, f_2: A_2 \to A$ are maps in
$\mathcal{C}_\Lambda$, $x \in U(A)$, $x_1 \in U(A_1)$, $x_2 \in U(A_2)$,
and $a_1, a_2 \in R(A)$ with $s(a_1) = U(f_1)(x_1)$,
$t(a_1) = x$ and $s(a_2) = U(f_2)(x_2)$, $t(a_2) = x$.
We construct a fiber product lying over $A_1 \times_A A_2$
for this diagram in $[U/R]$ as follows.

\medskip\noindent
Let $a = c(i(a_1), a_2)$, where $i: R \to R$ is the inverse morphism.
Then $a \in R(A)$, $x_2 \in U(A_2)$ and $s(a) = U(f_2)(x_2)$.
Hence an element $(a, x_2) \in R(A) \times_{s, U(A), U(f_2)} U(A_2)$.
By smoothness of $s : R \to U$ there is an element
$\widetilde{a} \in R(A_2)$ with $R(f_2)(\widetilde{a}) = a$ and
$s(\widetilde{a}) = x_2$. In particular
$U(f_2)(t(\widetilde{a})) = t(a) = U(f_1)(x_1)$. Thus $x_1$ and
$t(\widetilde{a})$ define an element
$$
(x_1, t(\widetilde{a})) \in U(A_1) \times_{U(A)} U(A_2).
$$
By the assumption that $U$ satisfies (RS), we have an identification
$U(A_1) \times_{U(A)} U(A_2) = U(A_1 \times_A A_2)$. Let us denote
$x_1 \times t(\widetilde{a}) \in U(A_1 \times_A A_2)$ the element
corresponding to $(x_1, t(\widetilde{a})) \in U(A_1) \times_{U(A)} U(A_2)$.
Let $p_1, p_2$ be the projections of $A_1 \times_A A_2$.  We claim
$$
\xymatrix{
(A_1 \times_A A_2, x_1 \times t(\widetilde{a})) \ar[d]_{(p_1, e(x_1))}
\ar[rr]_-{(p_2, i(\widetilde{a}))} & & (A_2, x_2) \ar[d]^{(f_2, a_2)} \\
(A_1, x_1) \ar[rr]^{(f_1, a_1)} & & (A, x)
}
$$
is a fiber square in $[U/R]$. (Note $e: U \to R$ denotes the identity.)

\medskip\noindent
The diagram is commutative because
$c(a_2, R(f_2)(i(\widetilde{a}))) = c(a_2, i(a)) = a_1$.
To check it is a fiber square, let
$$
\xymatrix{
(B, z) \ar[d]_{(g_1, b_1)} \ar[rr]_{(g_2, b_2)} & & (A_2, x_2)
\ar[d]^{(f_2, a_2)} \\
(A_1, x_1) \ar[rr]^{(f_1, a_1)} & & (A, x)
}
$$
be a commutative diagram in $[U/R]$. We will show there is a unique morphism
$(g, b) : (B, z) \to (A_1 \times_A A_2, x_1 \times t(\widetilde{a}))$
compatible with the morphisms to $(A_1, x_1)$ and $(A_2, x_2)$.
We must take $g = (g_1, g_2) : B \to A_1 \times_A A_2$.
Since by assumption $R$ satisfies (RS), we have an identification
$R(A_1 \times_A A_2) = R(A_1) \times_{R(A)} R(A_2)$.
Hence we can write $b = (b'_1, b'_2)$ for some
$b'_1 \in R(A_1)$, $b'_2 \in R(A_2)$ which agree in $R(A)$.
Then
$((g_1, g_2), (b'_1, b'_2)) : (B, z) \to
(A_1 \times_A A_2, x_1 \times t(\widetilde{a}))$
will commute with the projections if and only if
$b'_1 = b_1$ and $b'_2 = c(\widetilde{a}, b_2)$ proving unicity and
existence.
\end{proof}

\end{document}
